%============================================================
% TEMPLATE SKRIPSI — UNIVERSITAS AMIKOM PURWOKERTO
% Engine  : XeLaTeX (font Times New Roman dari sistem)
% Compile : latexmk -xelatex main.tex   (atau: make skripsi)
% Overleaf: set Menu > Compiler > XeLaTeX
%============================================================
% Satu file konten diedit sekali, seluruh dokumen tersusun dari:
%   config/    — identitas, toggle, paket, font, format
%   formats/   — halaman-halaman awal (cover s.d. motto)
%   pages/     — kata pengantar, daftar-daftar, abstrak, lampiran
%   chapters/  — Bab I s.d. Bab V
%   refs/      — daftar pustaka (BibTeX)
%============================================================
\documentclass[12pt,a4paper,oneside]{report}

%============================================================
% VARIABEL IDENTITAS DOKUMEN — SKRIPSI UNIVERSITAS AMIKOM PURWOKERTO
% Edit bagian ini sesuai dengan data skripsi Anda.
% File ini di-input paling awal agar semua file lain dapat
% mereferensikan variabel-variabel berikut.
%============================================================

% ── Identitas Mahasiswa ─────────────────────────────────────
% Nama: Capitalize Each Word (huruf awal tiap kata kapital).
% NIM:  ditulis dengan format titik sesuai kampus, contoh 20.21.0001.
\newcommand{\NamaMahasiswa}{Nama Mahasiswa}
\newcommand{\NIM}{xx.xx.xxxx}

% ── Judul Skripsi ────────────────────────────────────────────
% Ditulis HURUF BESAR dan tidak disingkat, maksimal 20 kata.
% Jika lebih dari 1 baris, atur pergantian baris dengan `\\`
% agar membentuk piramida terbalik (huruf V).
% Uppercase diterapkan otomatis di tempat yang membutuhkannya.
\newcommand{\JudulSkripsi}{Judul Skripsi Ditulis Huruf Besar\\
Dan Tidak Disingkat, Maksimal Dua Puluh Kata}

% Studi kasus (opsional). Kosongkan jika tidak ada.
% Contoh: (Studi Kasus: PT. Sinergi Dimensi Informatika)
\newcommand{\JudulStudiKasus}{(Studi Kasus: Jika Ada)}

% ── Dosen Pembimbing (gelar terbaru, NIDN tidak disingkat) ──
\newcommand{\NamaPembimbingSatu}{Nama Dosen Pembimbing 1, S.Kom., M.Kom.}
\newcommand{\NIDNPembimbingSatu}{00xx0xxxxxxxx}
\newcommand{\NamaPembimbingDua}{Nama Dosen Pembimbing 2, S.Kom., M.Kom.}
\newcommand{\NIDNPembimbingDua}{00xx0xxxxxxxx}

% ── Dosen Penguji (gelar terbaru, NIDN tidak disingkat) ─────
\newcommand{\NamaPengujiSatu}{Nama Penguji 1, S.Kom., M.Kom.}
\newcommand{\NIDNPengujiSatu}{00xx0xxxxxxxx}
\newcommand{\NamaPengujiDua}{Nama Penguji 2, S.Kom., M.Kom.}
\newcommand{\NIDNPengujiDua}{00xx0xxxxxxxx}

% Penguji 3 / Pendamping — diaktifkan lewat config/toggles.tex
\newcommand{\NamaPengujiTiga}{Nama Penguji 3 / Pendamping, S.Kom., M.Kom.}
\newcommand{\NIDNPengujiTiga}{00xx0xxxxxxxx}

% ── Ketua Program Studi (gelar + NIK) ────────────────────────
\newcommand{\NamaKaprodi}{Nama Ketua Program Studi, S.Kom., M.Kom.}
\newcommand{\NIKKaprodi}{00xx0xxxxxxxx}

% ── Dekan Fakultas ───────────────────────────────────────────
% Default di bawah adalah Dekan Fakultas Ilmu Komputer.
% Sesuaikan bila Fakultas Bisnis dan Ilmu Sosial.
\newcommand{\NamaDekan}{Prof. Dr. Eng. Ir. Imam Tahyudin, M.M.}
\newcommand{\NIKDekan}{2012.09.1.009}

% ── Program Studi & Fakultas ─────────────────────────────────
% Fakultas: "Fakultas Ilmu Komputer" (Prodi SI, TI, Informatika)
%           atau "Fakultas Bisnis dan Ilmu Sosial" (Bisnis Digital,
%           Ilmu Komunikasi). Uppercase diterapkan otomatis di cover.
\newcommand{\NamaProdi}{Informatika}
\newcommand{\NamaFakultas}{Fakultas Ilmu Komputer}

% ── Gelar yang Diperoleh ─────────────────────────────────────
\newcommand{\GelarSkripsi}{Sarjana Komputer (S.Kom.)}

% ── Tahun & Tanggal ──────────────────────────────────────────
\newcommand{\TahunUjian}{2026}            % tahun ujian skripsi
\newcommand{\TanggalACC}{1 Januari 2026}  % tanggal ACC terakhir dosen pembimbing
\newcommand{\TanggalUjian}{20 Januari 2026} % tanggal ujian/sidang skripsi
\newcommand{\TanggalKataPengantar}{Purwokerto, \today}  % penulisan kata pengantar
\newcommand{\KotaPernyataan}{Purwokerto, \TanggalUjian} % tempat/tanggal surat pernyataan

% ── Kata Kunci Abstrak (3–5 kata) ────────────────────────────
% Default di bawah = placeholder template asli (Template_Skripsi.docx).
\newcommand{\KataKunci}{satu, dua, tiga, empat, lima}
\newcommand{\Keywords}{satu, dua, tiga, empat, lima}
   % identitas skripsi (edit dulu!)
%============================================================
% TOGGLE KONFIGURASI
% Ubah nilai true/false sesuai kondisi skripsi Anda.
%============================================================

% Dua pembimbing atau pembimbing tunggal?
%   true  -> Halaman Persetujuan memuat Dosen Pembimbing 1 & 2
%   false -> Halaman Persetujuan memuat Pembimbing Tunggal saja
\newif\ifduapembimbing
\duapembimbingtrue

% Tampilkan Penguji 3 / Pendamping di Halaman Pengesahan?
\newif\ifpengujitiga
\pengujitigatrue

% Tampilkan halaman Daftar Istilah? (template asli: "jika ada")
\newif\ifdaftaristilah
\daftaristilahfalse
     % pilihan halaman (2 pembimbing, dll.)
%============================================================
% PAKET-PAKET LATEX (khusus engine XeLaTeX)
% Compile dengan: latexmk -xelatex main.tex
% Font Times New Roman diambil dari font sistem
% (tersedia di Windows/macOS dan Overleaf).
%============================================================

% ── Font & dasar ─────────────────────────────────────────────
\usepackage{fontspec}                 % pengganti inputenc/fontenc di XeLaTeX
\setmainfont{Times New Roman}         % sesuai ketentuan kampus
\usepackage[a4paper,
            left=4cm, right=3cm,      % penjilidan: kiri 4 cm
            top=4cm,  bottom=3cm]{geometry}
\usepackage{setspace}                 % \singlespacing, \onehalfspacing, \doublespacing
\usepackage{indentfirst}              % paragraf pertama juga menjorok
\usepackage{graphicx}
\usepackage{xcolor}                   % dipakai untuk teks tersembunyi (putih)
\usepackage{array}
\usepackage{longtable}                % tabel multi-halaman dengan "(Lanjutan)"
\usepackage{enumitem}
\usepackage{amsmath}                  % rumus matematika + \eqref

% ── Judul bab/sub-bab & daftar isi ───────────────────────────
\usepackage{titlesec}
\usepackage[titles]{tocloft}          % opsi `titles`: biar tidak bentrok titlesec

% ── Caption tabel & gambar ───────────────────────────────────
\usepackage{caption}

% ── Header/footer & halaman ──────────────────────────────────
\usepackage{fancyhdr}
\usepackage{pdfpages}                 % memasukkan scan (kartu bimbingan) di lampiran
\usepackage{etoolbox}

% ── Bahasa & bibliografi ─────────────────────────────────────
\usepackage[indonesian]{babel}
\usepackage{csquotes}
\usepackage[%
  backend=biber,
  style=authoryear,
  sorting=nyt,          % urut nama (belakang dulu), tahun, judul — alfabetis
  giveninits=true,      % nama depan ditulis inisial: Wibowo, A. R.
  maxbibnames=2,        % 1–2 pengarang di daftar pustaka: tulis semua;
  minbibnames=1,        % 3+ pengarang : pengarang pertama + "dkk."
  maxnames=2,           % batas yang sama untuk sitasi dalam teks
  minnames=1,
  uniquename=false,
  uniquelist=false,
  dashed=false          % pengarang berulang tetap ditulis penuh
]{biblatex}
\addbibresource{refs/references.bib}

% ── Hyperlink (selalu terakhir) ──────────────────────────────
\usepackage[hidelinks]{hyperref}
    % paket + biblatex
%============================================================
% MAKRO FONT
% Seluruh naskah memakai Times New Roman 12pt (default kelas).
% Makro berikut hanya untuk peran yang berbeda ukuran.
%============================================================

% Isi tabel: Times New Roman 10, spasi 1, rata kanan-kiri (justify)
\newcommand{\FontTabel}{\footnotesize}

% Judul bab & sub-bab: Times New Roman 12, bold
\newcommand{\FontJudulBab}{\normalsize\bfseries}
\newcommand{\FontJudulSubBab}{\normalsize\bfseries}
       % makro font
%============================================================
% FORMAT FORMAT AKHIR — Universitas Amikom Purwokerto
% Semua aturan format naskah dikumpulkan di sini:
%   - paragraf & spasi
%   - format BAB / sub-bab (A., B., ...)
%   - daftar isi / tabel / gambar / lampiran
%   - nomor halaman
%   - caption tabel & gambar
%   - gaya Daftar Pustaka Amikom (biblatex)
%   - makro-makro bantu template
%============================================================

% ── Paragraf ─────────────────────────────────────────────────
\setlength{\parindent}{1.27cm}    % perataan awal paragraf (standar Word 0.5")
\setlength{\parskip}{0pt}         % tanpa jarak antar-paragraf
\clubpenalty=10000                % cegah baris yatim/piatu
\widowpenalty=10000

% Spasi utama naskah (Bab I–V, termasuk judul bab): spasi 2.
% Halaman-halaman awal mengatur spasinya sendiri di file masing-masing.
\doublespacing

% ── Penomoran ────────────────────────────────────────────────
% Bab di daftar isi & judul pakai angka Romawi: BAB I, BAB II, ...
% Tabel/gambar/rumus tetap arab: Tabel 2.1, Gambar 2.1, Persamaan (2.1).
\renewcommand{\thechapter}{\Roman{chapter}}
\renewcommand{\thetable}{\arabic{chapter}.\arabic{table}}
\renewcommand{\thefigure}{\arabic{chapter}.\arabic{figure}}
\renewcommand{\theequation}{\arabic{chapter}.\arabic{equation}}
% Sub-bab pakai huruf: A., B., C. (bukan 1.1, 1.2) — sesuai template asli.
\renewcommand{\thesection}{\Alph{section}}
\renewcommand{\thesubsection}{\arabic{subsection}}
\renewcommand{\thesubsubsection}{\alph{subsubsection}}

% ── Format judul bab ─────────────────────────────────────────
% Dua baris terpusat bold:
%       BAB I
%       PENDAHULUAN
% Spasi 2 tetap berlaku di judul bab (termasuk judul bab).
\titleformat{\chapter}[display]
  {\centering\normalsize\bfseries}
  {BAB \thechapter}
  {0.6em}
  {\MakeUppercase}
% Sebelum/sesudah judul bab dikontrol ketat agar tidak ada jarak berlebih.
\titlespacing*{\chapter}{0pt}{-1.2em}{0.6em}

% Sub-bab:  A. Latar Belakang Masalah   (bold, rata kiri)
\titleformat{\section}
  {\normalsize\bfseries}{\thesection.}{0.5em}{}
\titlespacing*{\section}{0pt}{0.4em}{0.2em}

% Sub-sub-bab:  1. Judul   (bold, rata kiri)
\titleformat{\subsection}
  {\normalsize\bfseries}{\thesubsection.}{0.5em}{}
\titlespacing*{\subsection}{0pt}{0.4em}{0.2em}

% Tingkat keempat:  a. Judul   (bold, rata kiri)
\titleformat{\subsubsection}
  {\normalsize\bfseries}{\thesubsubsection.}{0.5em}{}
\titlespacing*{\subsubsection}{0pt}{0.4em}{0.2em}

% ── Daftar isi / tabel / gambar (tocloft) ────────────────────
% Hanya BAB dan sub-bab yang masuk daftar isi (level 2), dan
% keduanya dicetak tebal. Teks daftar: Times New Roman 12,
% spasi 1,5, tidak di-bold (kecuali bab & sub-bab).
\setcounter{tocdepth}{1}
\setcounter{secnumdepth}{3}       % sub-sub-bab bernomor di naskah

\renewcommand{\cftchapfont}{\bfseries}
\renewcommand{\cftchappagefont}{\bfseries}
\renewcommand{\cftchappresnum}{BAB }
\renewcommand{\cftchapaftersnum}{\hspace{0.4em}}
\setlength{\cftchapnumwidth}{3.2em}        % cukup untuk "BAB IV"
\renewcommand{\cftchapleader}{\cftdotfill{\cftdotsep}}

\renewcommand{\cftsecfont}{\bfseries}
\renewcommand{\cftsecpagefont}{\bfseries}
\renewcommand{\cftsecfont}{\bfseries}
\renewcommand{\cftsecpagefont}{\bfseries}
\renewcommand{\cftsecaftersnum}{.}
\setlength{\cftsecnumwidth}{1.8em}         % lebar "A."
\setlength{\cftsecindent}{1.5em}
\setlength{\cftbeforechapskip}{0em}        % rapat agar daftar isi muat satu halaman
% Spasi judul halaman daftar (DAFTAR ISI/TABEL/GAMBAR/LAMPIRAN)
% dirapatkan agar daftar isi tidak meluber ke halaman berikutnya.
\setlength{\cftbeforetoctitleskip}{1em}
\setlength{\cftaftertoctitleskip}{1em}

% Entri DAFTAR TABEL / DAFTAR GAMBAR memakai prefix seperti
% caption-nya: "Tabel 2.1. Data mahasiswa" / "Gambar 2.1. Logo lama".
\renewcommand{\cfttabpresnum}{Tabel~}
\renewcommand{\cfttabaftersnum}{.}
\setlength{\cfttabnumwidth}{4.6em}
\renewcommand{\cftfigpresnum}{Gambar~}
\renewcommand{\cftfigaftersnum}{.}
\setlength{\cftfignumwidth}{4.6em}

% ── Caption tabel & gambar ───────────────────────────────────
% Tabel  : caption DI ATAS tabel, format "Tabel X.Y. Judul".
% Gambar : caption DI BAWAH gambar, format "Gambar X.Y. Judul".
\captionsetup{%
  labelsep=period,
  font={normalsize,stretch=1},
  justification=centering,
  singlelinecheck=true}
\captionsetup[table]{position=above,skip=6pt}
\captionsetup[figure]{position=below,skip=6pt}

% ── Nomor halaman ────────────────────────────────────────────
% Awal halaman bab & halaman bagian awal : bawah tengah (plain)
% Halaman lanjutan                        : kanan atas  (utama)
% DAFTAR PUSTAKA & LAMPIRAN               : tanpa nomor halaman
\setlength{\headheight}{14.5pt}       % cukup untuk fancyhdr
\fancypagestyle{plain}{%
  \fancyhf{}%
  \fancyfoot[C]{\thepage}%
  \renewcommand{\headrulewidth}{0pt}%
  \renewcommand{\footrulewidth}{0pt}%
}
\fancypagestyle{utama}{%
  \fancyhf{}%
  \fancyhead[R]{\thepage}%
  \renewcommand{\headrulewidth}{0pt}%
  \renewcommand{\footrulewidth}{0pt}%
}

% ── Nama daftar ──────────────────────────────────────────────
% babel-indonesian menimpa nama-nama ini saat \begin{document},
% jadi perubahan dibungkus \AtBeginDocument.
\AtBeginDocument{%
  \renewcommand{\contentsname}{DAFTAR ISI}%
  \renewcommand{\listtablename}{DAFTAR TABEL}%
  \renewcommand{\listfigurename}{DAFTAR GAMBAR}%
}

%============================================================
% MAKRO-MAKRO BANTU TEMPLATE
%============================================================

% \TeksTersembunyi{teks} — teks putih (tak terlihat di PDF).
% Template asli menaruh tulisan "HALAMAN SAMPUL", "HALAMAN JUDUL",
% dst. berwarna putih agar dapat terbaca oleh generator daftar isi.
\newcommand{\TeksTersembunyi}[1]{\textcolor{white}{#1}}

% \HalamanAwal{NAMA HALAMAN} — teks putih tersembunyi + masuk
% DAFTAR ISI (tanpa nomor bab, dengan nomor halaman romawi).
\makeatletter
\newcommand{\HalamanAwal}[1]{%
  \par\noindent\TeksTersembunyi{#1}\par
  \addcontentsline{toc}{halamanawal}{#1}%
}
% Entri daftar isi untuk halaman awal: normal (tidak tebal), ada titik-titik
% dan nomor halaman romawi.
\newcommand{\l@halamanawal}{\@dottedtocline{1}{0pt}{0pt}}

% \TocHalaman{NAMA} — entri DAFTAR ISI untuk halaman yang sudah
% punya judul terlihat (KATA PENGANTAR, DAFTAR ISI, DAFTAR TABEL,
% ABSTRAK, dan sebagainya). Tanpa teks putih tambahan.
\newcommand{\TocHalaman}[1]{\addcontentsline{toc}{halamanawal}{#1}}

% \BagianAkhir{NAMA BAGIAN} — entri DAFTAR PUSTAKA / LAMPIRAN:
% tebal, TANPA titik-titik dan TANPA nomor halaman (sesuai template asli).
\newcommand{\BagianAkhir}[1]{%
  \addtocontents{toc}{\protect\contentsline{bagianakhir}{\protect#1}{}{}}%
}
\newcommand{\l@bagianakhir}[2]{%
  \vskip 0.1em
  \noindent\textbf{#1}\par
}
\makeatother

% \Istilah{istilah}{penjelasan} — entri DAFTAR ISTILAH.
\newcommand{\Istilah}[2]{%
  \noindent\textbf{#1} : #2\par
}

% \Lampiran{Judul} — mulai lampiran baru: "Lampiran N. Judul"
% (bold), otomatis masuk DAFTAR LAMPIRAN.
\newcounter{lampiran}
\newcommand{\LampiranList}[1]{\noindent #1\par\vspace{0.4em}}
\makeatletter
\newcommand{\Lampiran}[1]{%
  \clearpage
  \stepcounter{lampiran}%
  \noindent\textbf{Lampiran \arabic{lampiran}. #1}\par\vspace{0.5em}
  \addtocontents{lampl}{\protect\LampiranList{Lampiran \arabic{lampiran}. #1}}%
}
% \listoflampiran — halaman DAFTAR LAMPIRAN (tanpa titik-titik
% dan tanpa nomor halaman pada entri-entrinya).
\newcommand{\listoflampiran}{%
  \chapter*{DAFTAR LAMPIRAN}%
  \TocHalaman{DAFTAR LAMPIRAN}%
  \@starttoc{lampl}%
}
\makeatother

% \KotakMaterai — kotak pengganti materai Rp10.000 pada Surat
% Pernyataan Keaslian Penelitian (diganti materai asli saat
% dicetak dan ditandatangani).
\newcommand{\KotakMaterai}{%
  \fbox{\parbox[c][2.8cm][c]{4.5cm}{\centering Bermaterai\\ Rp10.000}}%
}

% \KotakScan — kotak penanda tempat memasukkan hasil scan pada
% lampiran (kartu bimbingan, surat keterangan, dll.).
\newcommand{\KotakScan}[1]{%
  \noindent\fbox{\parbox[c][8cm][c]{\dimexpr\linewidth-2\fboxsep-2\fboxrule\relax}{%
    \centering #1}}%
}

% ── Lingkungan tabel (longtable) ─────────────────────────────
% \begin{tabel}{spesifikasi-kolom} ... \end{tabel}
% Isi tabel otomatis: Times New Roman 10, spasi 1, justify.
% Lihat contoh pemakaian lengkap (dengan tabel lanjutan) di
% chapters/bab2.tex.
\newenvironment{tabel}[1]{%
  \begingroup
  \FontTabel
  \setstretch{1.0}
  \setlength{\LTpre}{6pt}%
  \setlength{\LTpost}{6pt}%
  \begin{longtable}{#1}%
}{%
  \end{longtable}%
  \endgroup
}

% \TabelLanjutan{jumlah-kolom}{Judul Tabel} — kepala halaman
% lanjutan tabel: "Tabel X.Y. Judul Tabel (Lanjutan)" + kepala tabel.
% Pemakaian di dalam \endhead (lihat contoh chapters/bab2.tex).
\newcommand{\TabelLanjutan}[2]{%
  \multicolumn{#1}{c}{Tabel \thetable. #2 (Lanjutan)}\\
}

% ── Halaman landscape (kertas miring) ────────────────────────
% Nomor halaman pada halaman miring tetap terbaca horizontal di
% KANAN ATAS orientasi miring (sesuai contoh halaman landscape di
% template asli). Caranya: glyph dicetak di kiri-atas lembar fisik
% dengan rotasi 90°, sehingga setelah halaman dirotasi viewer
% (/Rotate) ia tampil horizontal di pojok kanan atas.
\fancypagestyle{landskape}{%
  \fancyhf{}%
  \fancyhead[L]{\rotatebox{90}{\thepage}}%
  \renewcommand{\headrulewidth}{0pt}%
  \renewcommand{\footrulewidth}{0pt}%
}
% Lingkungan landskape — pembungkus nyaman:
%   \begin{landskape} ... konten/tabel ... \end{landskape}
\newenvironment{landskape}{%
  \begin{landscape}%
  \pagestyle{landskape}%
}{%
  \end{landscape}%
  \pagestyle{utama}%
}

%============================================================
% DAFTAR PUSTAKA — GAYA AMIKOM (biblatex)
%============================================================
% Judul ditulis apa adanya (Title Case dari file .bib, tidak
% dirubah jadi sentence case oleh biblatex).
\DeclareFieldFormat{titlecase}{#1}

% Pengarang terakhir dihubungkan "dan" (bukan "&"):
%   Wibowo, A. R., dan Saputra, I. A.
\DeclareDelimFormat[bib,biblist]{finalnamedelim}{\addcomma\space dan}

% Kata-kata lokal ("dan", "dkk.", "Diambil dari", dsb.) disediakan
% oleh indonesian.lbx — file lokalisasi biblatex milik repo,
% karena biblatex belum menyediakan indonesian.lbx.
% Judul artikel/tesis tanpa tanda kutip (gaya Amikom: judul polos,
% yang italic hanya nama jurnal/judul buku).
\DeclareFieldFormat[article,inbook,incollection,inproceedings,patent,thesis,unpublished]{title}{#1\isdot}
% Volume(nomor) halaman tanpa awalan "pp." : "Jurnal Teknik Elektro, 11(1), 26-30"
\DeclareFieldFormat{pages}{#1}
\renewbibmacro*{journal+issuetitle}{%
  \usebibmacro{journal}%
  \setunit*{\addspace}%
  \printfield[journaltitle]{series}%
  \setunit{\addcomma\space}%
  \printfield[nq]{issue}%
  \setunit{\addcomma\space}%
  \usebibmacro{volume+number+eid}%
  \setunit{\addcomma\space}%
  \usebibmacro{issue+date}%
  \setunit{\addcolon\space}%
  \usebibmacro{issue}%
  \newunit}
\renewbibmacro*{volume+number+eid}{%
  \printfield{volume}%
  \iffieldundef{number}{}{\mkbibparens{\printfield{number}}}%
  \setunit{\addcomma\space}%
  \printfield{eid}%
}
% Jenis skripsi: "Skripsi." (bukan "Master's thesis.")
\DeclareFieldFormat[thesis]{type}{#1\adddot}

% Hilangkan awalan "in:" pada prosiding (gaya Amikom tidak memakainya)
\renewbibmacro{in:}{}

% Jarak penulisan: dalam referensi spasi 1, antar referensi 2 spasi.
\renewcommand{\bibfont}{\setstretch{1.0}}
\setlength{\bibitemsep}{2\baselineskip}
\setlength{\bibhang}{1.27cm}
  % seluruh aturan format

\begin{document}

% ── BAGIAN AWAL (nomor halaman romawi, mulai i di sampul) ──
\pagestyle{plain}            % nomor halaman: bawah tengah
\pagenumbering{roman}

%============================================================
% HALAMAN SAMPUL DEPAN (COVER)
%   - Judul: uppercase bold, piramida terbalik (atur dengan \\)
%   - Teks putih "HALAMAN SAMPUL" di antara judul dan "Skripsi"
%   - Logo: tinggi 4 cm, lebar 5,5 cm, terpusat
%   - Identitas: nama 12pt bold, NIM, prodi, fakultas, kampus, tahun
% Halaman ini dihitung halaman i (romawi) tetapi nomornya tidak
% ditampilkan.
%============================================================
\thispagestyle{empty}
\onehalfspacing

\begin{center}
  % ── Judul skripsi (piramida terbalik) ──
  \vspace*{1em}
  {\bfseries\MakeUppercase{\JudulSkripsi}\par}
  \JudulStudiKasus

  % ── Teks tersembunyi putih: "HALAMAN SAMPUL" ──
  \HalamanAwal{HALAMAN SAMPUL}

  \vspace{0.4em}
  {\bfseries Skripsi\par}

  % ── Logo kampus: tinggi 4 cm, lebar maksimal 5,5 cm ──
  \vfill
  \includegraphics[height=4cm,width=5.5cm,keepaspectratio]{assets/amikom/logo_amikom.png}

  % ── Identitas mahasiswa ──
  \vfill
  Disusun oleh

  \vspace{0.3em}
  {\bfseries \NamaMahasiswa\par}
  {\bfseries \NIM\par}

  % ── Program studi, fakultas, universitas, kota, tahun ujian ──
  \vfill
  {\bfseries PROGRAM STUDI \MakeUppercase{\NamaProdi}\par}
  {\bfseries \MakeUppercase{\NamaFakultas}\par}
  {\bfseries UNIVERSITAS AMIKOM PURWOKERTO\par}
  {\bfseries PURWOKERTO\par}
  \vspace{0.4em}
  {\bfseries \TahunUjian\par}
\end{center}

\clearpage
          % HALAMAN SAMPUL        (i)
%============================================================
% HALAMAN JUDUL DALAM
%   Sama seperti sampul depan, ditambah:
%   - teks putih "HALAMAN JUDUL"
%   - kalimat persyaratan derajat (spasi 1)
%============================================================
\setstretch{1.43}

\begin{center}
  % ── Judul skripsi (piramida terbalik) ──
  \vspace*{1em}
  {\bfseries\MakeUppercase{\JudulSkripsi}\par}
  \JudulStudiKasus

  % ── Teks tersembunyi putih: "HALAMAN JUDUL" ──
  \HalamanAwal{HALAMAN JUDUL}

  \vspace{0.4em}
  {\bfseries Skripsi\par}

  % ── Kalimat persyaratan (spasi 1) ──
  \begin{spacing}{1.0}
    Untuk memenuhi sebagian persyaratan mencapai derajat Sarjana S1\\
    pada Program Studi \NamaProdi\ \NamaFakultas
  \end{spacing}

  % ── Logo kampus: tinggi 4 cm, lebar maksimal 5,5 cm ──
  \vfill
  \includegraphics[height=4cm,width=5.5cm,keepaspectratio]{assets/amikom/logo_amikom.png}

  % ── Identitas mahasiswa ──
  \vfill
  Disusun oleh

  \vspace{0.3em}
  {\bfseries \NamaMahasiswa\par}
  {\bfseries \NIM\par}

  % ── Program studi, fakultas, universitas, kota, tahun ujian ──
  \vfill
  {\bfseries PROGRAM STUDI \MakeUppercase{\NamaProdi}\par}
  {\bfseries \MakeUppercase{\NamaFakultas}\par}
  {\bfseries UNIVERSITAS AMIKOM PURWOKERTO\par}
  {\bfseries PURWOKERTO\par}
  \vspace{0.4em}
  {\bfseries \TahunUjian\par}
\end{center}

\clearpage
  % HALAMAN JUDUL
%============================================================
% HALAMAN PERSETUJUAN DOSEN PEMBIMBING
%   Mendukung 2 Pembimbing atau Pembimbing Tunggal
%   (atur di config/toggles.tex: \duapembimbingtrue / \duapembimbingfalse)
%   Isi: judul, mahasiswa, tanggal ACC terakhir, blok tanda tangan
%   pembimbing (gelar terbaru + NIDN penuh), pernyataan "telah
%   diterima", serta mengetahui Dekan dan Kaprodi.
%============================================================
\chapter*{PERSETUJUAN}
\HalamanAwal{HALAMAN PERSETUJUAN}

\begin{center}
  {\bfseries Skripsi\par}
  \vspace{0.4em}
  {\bfseries\MakeUppercase{\JudulSkripsi}\par}
  \JudulStudiKasus

  \vspace{1.2em}
  Yang dipersiapkan dan disusun oleh

  \vspace{0.3em}
  {\bfseries \NamaMahasiswa\par}
  {\bfseries NIM. \NIM\par}
\end{center}

\vspace{0.4em}
telah disetujui oleh dosen pembimbing Skripsi\\
Pada tanggal \TanggalACC

% ── Blok tanda tangan dosen pembimbing (spasi 1) ──
\begin{spacing}{1.0}
\ifduapembimbing
  \noindent
  \begin{minipage}[t]{0.48\linewidth}
    \centering
    Dosen Pembimbing 1,\\
    \vspace{2.2cm}
    \textbf{\NamaPembimbingSatu}\\
    NIDN. \NIDNPembimbingSatu
  \end{minipage}%
  \hfill
  \begin{minipage}[t]{0.48\linewidth}
    \centering
    Dosen Pembimbing 2,\\
    \vspace{2.2cm}
    \textbf{\NamaPembimbingDua}\\
    NIDN. \NIDNPembimbingDua
  \end{minipage}
\else
  \begin{center}
    Dosen Pembimbing 1,\\
    \vspace{2.2cm}
    \textbf{\NamaPembimbingSatu}\\
    NIDN. \NIDNPembimbingSatu
  \end{center}
\fi

\vspace{1em}
% ── Pernyataan telah diterima (spasi 1) ──
\begin{center}
  Skripsi ini telah diterima sebagai salah satu persyaratan\\
  untuk memperoleh gelar \GelarSkripsi\\
  Tanggal \TanggalUjian
\end{center}

\vspace{1em}
% ── Mengetahui: Dekan & Kaprodi ──
\noindent
\begin{minipage}[t]{0.48\linewidth}
  Mengetahui,\\
  \NamaFakultas\\
  Dekan,
  \vspace{2.2cm}\\
  \textbf{\NamaDekan}\\
  NIK. \NIKDekan
\end{minipage}%
\hfill
\begin{minipage}[t]{0.48\linewidth}
  Kaprodi \NamaProdi
  \vspace{2.2cm}\\
  \textbf{\NamaKaprodi}\\
  NIK. \NIKKaprodi
\end{minipage}
\end{spacing}

\clearpage
    % HALAMAN PERSETUJUAN
%============================================================
% HALAMAN PENGESAHAN DEWAN PENGUJI
%   Isi: tanggal ujian, tabel daftar penguji (beserta NIDN),
%   pernyataan pengesahan gelar S.Kom., dan tanda tangan Dekan.
%   Penguji 3 / Pendamping diatur di config/toggles.tex.
%   Blok penguji memakai tabular agar nama/NIDN antar kolom
%   sejajar pada baseline yang sama.
%============================================================
\chapter*{PENGESAHAN}
\HalamanAwal{HALAMAN PENGESAHAN}

\begin{center}
  {\bfseries Skripsi\par}
  \vspace{0.4em}
  {\bfseries\MakeUppercase{\JudulSkripsi}\par}
  \JudulStudiKasus

  \vspace{1.2em}
  Yang dipersiapkan dan disusun oleh

  \vspace{0.3em}
  {\bfseries \NamaMahasiswa\par}
  {\bfseries NIM. \NIM\par}
\end{center}

\vspace{1em}
\noindent Telah dipertahankan di depan dewan penguji\\
Pada tanggal \TanggalUjian

% ── Tabel daftar penguji (spasi 1) ──
\begin{spacing}{1.0}
\vspace{1.6em}
\ifpengujitiga
  \noindent\begin{tabular}{@{}>{\centering\arraybackslash}p{4.4cm}
                              >{\centering\arraybackslash}p{4.4cm}
                              >{\centering\arraybackslash}p{4.4cm}@{}}
    Penguji 1, & Penguji 2, & Penguji 3 / Pendamping,\\[2.0cm]
    \textbf{\NamaPengujiSatu} & \textbf{\NamaPengujiDua} &
    \textbf{\NamaPengujiTiga}\\
    NIDN. \NIDNPengujiSatu & NIDN. \NIDNPengujiDua &
    NIDN. \NIDNPengujiTiga\\
  \end{tabular}
\else
  \noindent\begin{tabular}{@{}>{\centering\arraybackslash}p{6.5cm}
                              >{\centering\arraybackslash}p{6.5cm}@{}}
    Penguji 1, & Penguji 2,\\[2.0cm]
    \textbf{\NamaPengujiSatu} & \textbf{\NamaPengujiDua}\\
    NIDN. \NIDNPengujiSatu & NIDN. \NIDNPengujiDua\\
  \end{tabular}
\fi

\vspace{1.2em}
% ── Pernyataan pengesahan (spasi 1) ──
\begin{center}
  Skripsi ini telah disahkan sebagai salah satu persyaratan\\
  untuk memperoleh gelar \GelarSkripsi\\
  Tanggal \TanggalUjian
\end{center}

\vspace{1em}
% ── Mengetahui Dekan ──
\noindent\begin{tabular}{@{}>{\raggedright\arraybackslash}p{7cm}@{}}
  Mengetahui,\newline \NamaFakultas\newline Dekan \\[2.2cm]
  \textbf{\NamaDekan}\\
  NIK. \NIKDekan\\
\end{tabular}
\end{spacing}

\clearpage
     % HALAMAN PENGESAHAN
%============================================================
% SURAT PERNYATAAN KEASLIAN PENELITIAN
%   Isi: identitas mahasiswa, pernyataan keaslian, blok tanda
%   tangan di KANAN-BAWAH (kota+tanggal, "Yang menyatakan,", kotak
%   materai Rp10.000, nama bergaris bawah, NIM) sesuai template
%   asli. Kotak "Bermaterai" adalah placeholder; materai asli
%   ditempel setelah naskah dicetak.
%   Catatan implementasi: isi baris tabular tidak boleh memuat
%   `\\` (memutus baris tabular) dan lingkungan spacing tidak
%   boleh membungkus tabular (menyebabkan baris tumpang tindih).
%============================================================
\chapter*{SURAT PERNYATAAN KEASLIAN PENELITIAN}
% Garis bawah ganda pada judul (border bawah ganda template asli),
% tepat di bawah judul — sebelum baris penanda putih
\vspace{-1.2em}
\noindent\rule{\linewidth}{0.5pt}\\[1.4pt]\rule{\linewidth}{0.5pt}
\HalamanAwal{HALAMAN PERNYATAAN KEASLIAN}

\begin{spacing}{1.43}
\noindent Yang bertanda tangan dibawah ini :
\end{spacing}

% ── Identitas mahasiswa (spasi 1,5 antar-baris seperti template
%    asli; dicapai lewat \arraystretch) ──
{\renewcommand{\arraystretch}{1.5}\noindent
\begin{tabular}{@{}p{4.2cm}@{:\ }l@{}}
  Nama Mahasiswa   & \NamaMahasiswa\\
  NIM              & \NIM\\
  Program Studi    & \NamaProdi\\
  Fakultas         & \NamaFakultas\\
  Perguruan Tinggi & Universitas Amikom Purwokerto\\
\end{tabular}\par}

\begin{spacing}{1.43}
\noindent Menyatakan bahwa Skripsi Sebagai Berikut:
\end{spacing}

% ── Judul & pembimbing (spasi 1,5) ──
{\renewcommand{\arraystretch}{1.5}
\noindent
\parbox[t]{4.2cm}{Judul Skripsi}\ :\ \parbox[t]{9.2cm}{\JudulSkripsi}

\medskip
\noindent\begin{tabular}{@{}p{4.2cm}@{:\ }l@{}}
  Dosen Pembimbing 1 & \NamaPembimbingSatu\\
  \ifduapembimbing
  Dosen Pembimbing 2 & \NamaPembimbingDua\\
  \else
  Dosen Pembimbing 2 & -\\
  \fi
\end{tabular}\par}

\begin{spacing}{1.43}
\noindent
Adalah benar-benar \textbf{ASLI} dan \textbf{BELUM PERNAH} dibuat
orang lain, kecuali yang diacu dalam daftar pustaka pada Skripsi
ini.

\noindent
Demikian pernyataan ini saya buat, apabila ini di kemudian hari
terbukti bahwa saya melakukan penjiplakan karya orang lain, maka
saya bersedia menerima \textbf{SANKSI AKADEMIK}.
\end{spacing}

% ── Blok tanda tangan: kanan-bawah, tersusun vertikal
%    (kota/tanggal → Yang menyatakan → materai → nama/NIM)
\begin{spacing}{1.0}
\vspace{0.2em}
\noindent\hspace*{5.8cm}%
\begin{minipage}[t]{7.4cm}
  \centering
  \KotaPernyataan

  \vspace{0.2em}
  Yang menyatakan,

  \vspace{0.4em}
  \KotakMaterai

  \vspace{1.0cm}
  \NamaTTD{\NamaMahasiswa}\\
  NIM. \NIM
\end{minipage}
\end{spacing}

\clearpage
     % HALAMAN PERNYATAAN KEASLIAN
%============================================================
% HALAMAN PERSEMBAHAN
%   Sesuai template asli: judul "HALAMAN PERSEMBAHAN" TERLIHAT
%   (bold, tengah) dan sekaligus masuk DAFTAR ISI. Isi spasi 2,
%   rata kiri-kanan, tanpa indentasi.
%   Teks di bawah disalin verbatim dari template asli
%   (reference/Template_Skripsi.docx) — ganti dengan persembahan
%   Anda.
%============================================================
\begin{center}
  {\bfseries HALAMAN PERSEMBAHAN\par}
\end{center}
\TocHalaman{HALAMAN PERSEMBAHAN}
\setstretch{1.9}

\noindent
Halaman ini berisi kepada siapa Skripsi dipersembahkan. Ditulis
dengan singkat, resmi, sederhana, tidak terlalu banyak, serta
tidak menjurus ke penulisan informal sehingga mengurangi sifat
resmi laporan ilmiah.

\noindent
Bagian ini berisi pernyataan resmi yang ingin disampaikan oleh
penulis kepada pihak lain, misalnya ucapan terima kasih kepada
Tim Dosen Pembimbing, Tim Dosen Penguji, dan semua pihak yang
terkait dalam penyelesaian Skripsi termasuk orang tua dan
penyandang dana.

\clearpage
    % HALAMAN PERSEMBAHAN
%============================================================
% HALAMAN MOTTO
%   Sesuai template asli: judul "HALAMAN MOTTO" TERLIHAT (bold,
%   tengah) dan sekaligus masuk DAFTAR ISI. Motto ditulis
%   singkat, resmi, sederhana (kata mutiara, ungkapan tokoh, atau
%   Kitab Suci); spasi 2.
%============================================================
\begin{center}
  {\bfseries HALAMAN MOTTO\par}
\end{center}
\TocHalaman{HALAMAN MOTTO}
\setstretch{1.9}

\begin{center}
  % Ganti dengan motto Anda (kata mutiara, ungkapan tokoh, atau Kitab Suci).
  \emph{``Sesungguhnya sesudah kesulitan itu ada kemudahan.''}

  \vspace{1em}
  ---  ---  ---  ---
\end{center}

\clearpage
          % HALAMAN MOTTO
%============================================================
% KATA PENGANTAR
%   Diakhiri kota, tanggal, dan kata "Penulis" rata kiri di
%   bawahnya (tanpa nama terang & tanda tangan) — sesuai template
%   asli. Spasi 2, paragraf menjorok 1,27cm.
%   Teks di bawah disalin verbatim dari template asli
%   (reference/Template_Skripsi.docx) — ganti dengan kata
%   pengantar Anda.
%============================================================
\chapter*{KATA PENGANTAR}
\TocHalaman{KATA PENGANTAR}
\setstretch{1.9}

Nama harus ditulis secara lengkap termasuk gelar akademik dan
harus dihindari ucapan terima kasih kepada pihak yang tidak
terkait. Bahasa yang digunakan harus mengikuti kaidah bahasa
Indonesia yang baku.

\noindent
Bagian ini tidak perlu dituliskan hal-hal yang bersifat ilmiah.
Kata Pengantar diakhiri dengan mencantumkan kota dan tanggal
penulisan diikuti di bawahnya dengan
\NamaTTD{kata ``Penulis'' tanpa perlu menyebutkan nama dan tanda
tangan}.

% Penutup: kota + tanggal, diikuti kata "Penulis" di bawahnya —
% keduanya RATA KIRI (sesuai template asli, tanpa nama terang dan
% tanpa tanda tangan).
\vspace{0.8em}
\noindent
\TanggalKataPengantar

\vspace{0.6em}
\noindent
Penulis

\clearpage
   % KATA PENGANTAR
%============================================================
% HALAMAN DAFTAR-DAFTAR
%   Daftar Isi → Daftar Tabel → Daftar Gambar →
%   Daftar Istilah (opsional) → Daftar Lampiran.
%   Teks daftar: Times New Roman 12, spasi 1,5, tidak tebal
%   (entri BAB & sub-bab otomatis tebal).
%============================================================

% ── DAFTAR ISI ──
{\onehalfspacing
\tableofcontents
\TocHalaman{DAFTAR ISI}
}

% ── DAFTAR TABEL ──
{\onehalfspacing
\listoftables
\TocHalaman{DAFTAR TABEL}
}

% ── DAFTAR GAMBAR ──
{\onehalfspacing
\listoffigures
\TocHalaman{DAFTAR GAMBAR}
}

% ── DAFTAR ISTILAH (jika ada — atur di config/toggles.tex) ──
\ifdaftaristilah
  %============================================================
% DAFTAR ISTILAH (jika ada)
%   Daftar istilah/istilah asing yang dipakai dalam skripsi
%   beserta penjelasannya. Ganti contoh di bawah dengan istilah
%   penelitian Anda.
%============================================================
\chapter*{DAFTAR ISTILAH}
\TocHalaman{DAFTAR ISTILAH}
\onehalfspacing

\Istilah{IoT}{Internet of Things, yaitu konsep penghubungan perangkat melalui internet untuk saling bertukar data.}
\Istilah{Load Cell}{Sensor pengukur berat/gaya yang mengubah besaran mekanik menjadi sinyal listrik.}
\Istilah{Mikrokontroler}{Integrated circuit yang dapat diprogram untuk mengendalikan perangkat elektronik.}

\fi

% ── DAFTAR LAMPIRAN (tanpa titik-titik & nomor halaman) ──
{\onehalfspacing
\listoflampiran
}
    % DAFTAR ISI/TABEL/GAMBAR/ISTILAH/LAMPIRAN
%============================================================
% ABSTRAK (Bahasa Indonesia)
%   1 paragraf, 150–250 kata, spasi 1, Times New Roman 12.
%   Diakhiri 3–5 kata kunci.
%   Teks di bawah disalin verbatim dari template asli
%   (reference/Template_Skripsi.docx) — ganti dengan abstrak Anda.
%   PENTING: \singlespacing dibungkus \begingroup...\endgroup —
%   tanpa itu deklarasinya bocor dan seluruh isi Bab I–V ikut
%   kehilangan spasi 2.
%============================================================
\begingroup
\singlespacing

\chapter*{ABSTRAK}
\TocHalaman{ABSTRAK}

\noindent
Abstrak merupakan \emph{outline} dari sebuah hasil
penelitian/karya ilmiah/naskah resmi yang memerlukan deskripsi
secara singkat. Abstrak disusun dengan kalimat yang singkat,
jelas, runtut, dan sistematis dan dapat menggambarkan isi laporan
secara keseluruhan. Abstrak disusun dalam bahasa Indonesia,
disusun menjadi 1 paragraf, berkisar antara 150-250 kata, diketik
dengan jarak 1 spasi, menggunakan jenis font Times New Roman,
size 12. Abstrak Skripsi memuat apa dan mengapa penelitian
dikerjakan, bagaimana dikerjakan, dan apa hasil penting yang
diperoleh dari penelitian. Abstrak berisi judul penelitian,
tujuan penelitian, batasan variabel penelitian, serta metode
analisis yang digunakan dalam penelitian. kemudian berisi
ringkasan hasil penelitian dan pembahasan. Selanjutnya,
mengungkapkan kesimpulan dari penelitian. Di bagian bawah
intisari dituliskan kata-kata kunci, bisa berupa kata-kata
penting dalam intisari atau kata yang sering muncul, berjumlah
minimal 3 (tiga) dan maksimal 5 (lima) kata.

\vspace{1em}
\noindent Kata kunci: \KataKunci.

\clearpage
\endgroup
          % ABSTRAK
%============================================================
% ABSTRACT (Bahasa Inggris)
%   Diterjemahkan dari abstrak bahasa Indonesia, dicetak miring
%   total, diakhiri Keywords (3–5 kata). Spasi 1.
%   Catatan: \itshape dibungkus group agar tidak bocor ke halaman
%   berikutnya.
%============================================================
% Judul "ABSTRACT" bold + ITALIC (seluruh halaman miring, sesuai
% template asli); entri DAFTAR ISI juga dicetak miring.
\chapter*{\itshape ABSTRACT}
\TocHalaman{\textit{ABSTRACT}}
\singlespacing

% Ganti teks panduan di bawah ini dengan terjemahan abstrak Anda.
% Seluruh isi abstrak ditulis dengan huruf miring (italic).
\begingroup
\itshape

\noindent
An abstract is an outline of the results of research that
requires a brief description. The abstract is composed of short,
clear, structured, and systematic sentences so that it can
describe the whole report. It states the research title, the
research objective, the research variable boundaries, and the
analysis method used, then summarizes the results and discussion,
and finally reveals the conclusion of the research.

\vspace{1em}
\noindent Keywords: \Keywords.

\endgroup

\clearpage
         % ABSTRACT

% ── BAGIAN INTI (nomor halaman arab, mulai 1 di Bab I) ──
\pagestyle{utama}            % halaman lanjutan: nomor kanan atas
\pagenumbering{arabic}
\setstretch{1.9}             % pastikan isi BAB spasi 2 — deklarasi
                             % spasi halaman awal tidak boleh bocor ke sini
% Awal setiap bab otomatis memakai nomor bawah tengah (plain).

%============================================================
% BAB I — PENDAHULUAN
% Struktur: A. Latar Belakang Masalah, B. Rumusan Masalah,
% C. Batasan Masalah, D. Tujuan Penelitian, E. Manfaat Penelitian.
% Spasi isi: 2 (termasuk judul bab). Sub-bab pakai huruf A., B., ...
%============================================================
\chapter{PENDAHULUAN}

\section{Latar Belakang Masalah}
% Ganti teks panduan di bawah ini dengan latar belakang Anda.
Bagian ini memuat penjelasan tentang fenomena umum yang terjadi
dan kemudian dipersempit mengarah ke permasalahan yang akan
diteliti atau dibahas. Latar belakang masalah merupakan uraian
yang komprehensif mengenai pentingnya permasalahan tersebut
diangkat menjadi topik penulisan, ditinjau dari aspek praktis
maupun kontribusi ilmiah secara teoritis
\parencite{winowatan2009}.

Penulisan latar belakang masalah disajikan dalam bentuk uraian
yang secara kronologis diarahkan untuk langsung menuju rumusan
masalah. Istilah asing ditulis miring, misalnya \emph{machine
learning} dan \emph{internet of things}. Semua simbol matematika,
variabel, dan persamaan wajib menggunakan sintaks mode matematika
LaTeX, misalnya $x$ dan $y$, atau tampilan rumus seperti pada
persamaan~\eqref{eq:contoh-bab2}.

\section{Rumusan Masalah}
Bagian ini memuat penjelasan tentang permasalahan sehingga
memerlukan solusi. Permasalahan yang diuraikan dalam latar
belakang masalah dirumuskan kembali secara tegas dan jelas dalam
bentuk poin-poin yang terinci yang berisi masalah-masalah yang
akan dikaji. Dirumuskan dalam bentuk pertanyaan atau pernyataan;
jumlahnya dapat satu atau lebih, tergantung kedalaman dan luasnya
masalah yang akan diteliti. Sebagai contoh:
\begin{enumerate}
  \item Bagaimana merancang dan membangun sistem yang diusulkan?
  \item Bagaimana menguji kelayakan sistem yang diusulkan?
\end{enumerate}

\section{Batasan Masalah}
Dalam penulisan batasan masalah harus dihindari penggunaan kata
``tidak dibahas/membahas''. Yang harus dituliskan adalah segala
sesuatu yang menjadi topik pembahasan skripsi. Misalnya:
``Penelitian ini dibatasi oleh hal-hal sebagai berikut:''
\begin{enumerate}
  \item Batasan-batasan permasalahan yang akan dicari solusinya;
  \item Asumsi-asumsi yang digunakan sebagai dasar penulisan;
  \item Batasan solusi.
\end{enumerate}

\section{Tujuan Penelitian}
Penelitian dapat bertujuan untuk mengetahui, menguraikan,
menerangkan, membuktikan, atau menerapkan model, konsep, atau
dugaan. Tujuan adalah kurang lebih dari jawaban/capaian yang ada
di rumusan masalah. Bagian ini memuat penjelasan secara spesifik:
\begin{enumerate}
  \item Solusi yang ditawarkan untuk menangani permasalahan yang dihadapi;
  \item Hal-hal yang ingin dicapai melalui kegiatan yang akan dilakukan.
\end{enumerate}

\section{Manfaat Penelitian}
Manfaat penelitian menguraikan dan menjelaskan kegunaan secara
teoritik dan aplikatif dari penelitian yang telah dilakukan.
Manfaat atau kontribusi penelitian memuat dua hal yang mendasar:
\begin{itemize}
  \item \textbf{Manfaat Teoritik}, apabila hasil penelitian akan
        menghasilkan sebuah pendapat baru atau hasil penerapan hukum;
  \item \textbf{Manfaat Aplikatif}, apabila terdapat manfaat atau
        nilai guna hasil penelitian bagi stakeholder atau pihak-pihak
        yang terkait langsung dengan hasil penelitian.
\end{itemize}

%============================================================
% BAB II — TINJAUAN PUSTAKA
% Struktur: A. Landasan Teori, B. Penelitian Sebelumnya.
% Bab ini sekaligus menjadi contoh pemakaian:
%   - tabel biasa (caption di atas)        : Tabel 2.1
%   - gambar  (caption di bawah)           : Gambar 2.1 & 2.2
%   - rumus matematika                     : Persamaan 2.1
%   - tabel multi-halaman "(Lanjutan)"     : Tabel 2.2
%   - tabel penelitian sebelumnya (3 artikel: 2 nasional,
%     1 internasional)                     : Tabel 2.3
%============================================================
\chapter{TINJAUAN PUSTAKA}

\section{Landasan Teori}
Menguraikan teori-teori yang menunjang penulisan/penelitian.
Referensi ataupun buku baru yang dirujuk harus dicantumkan
sebagai rujukan dan harus dicantumkan dalam daftar pustaka
\parencite{winowatan2009}. Perumusan landasan teori hendaknya
memenuhi ketentuan sebagai berikut:
\begin{enumerate}
  \item Pustaka yang digunakan terbaru, relevan, dan asli dari karya ilmiah;
  \item Apabila sumber informasi dan data yang dirujuk berasal dari buku,
        usahakan mencari terbitan edisi paling akhir;
  \item Penulisan landasan teori tidak mengandung unsur plagiat.
\end{enumerate}

Uraian teori yang memerlukan data pendukung dalam bentuk tabel
maupun gambar, format penulisan nama tabel dan gambar disajikan
seperti Tabel~\ref{tab:data-mahasiswa} dan
Gambar~\ref{fig:logo-lama}. Nomor urut tabel dan gambar
disesuaikan dengan nomor bab, diikuti dengan nomor urut dari tabel
atau gambar tersebut.

% ── Contoh tabel: caption di ATAS tabel, isi Times New Roman 10,
%    spasi 1, justify. Dipanggil di dalam paragraf (wajib). ──
\begin{tabel}{@{}c l l l@{}}
  \caption{Data mahasiswa}\label{tab:data-mahasiswa}\\
  \hline
  \rowcolor{gray!25}
  No & NIM & Nama Mahasiswa & Konsentrasi\\
  \hline
  \endfirsthead
  \TabelLanjutan{4}{Data mahasiswa}\\
  \hline
  \rowcolor{gray!25}
  No & NIM & Nama Mahasiswa & Konsentrasi\\
  \hline
  \endhead
  \hline
  \endfoot
  \endlastfoot
  1 & 11.11.0001 & Mahasiswa 1 & Konsentrasi 1\\
  2 & 11.11.0002 & Mahasiswa 2 & Konsentrasi 2\\
\end{tabel}

Terdapat perubahan bentuk logo Universitas Amikom Purwokerto,
logo sebelumnya ditunjukkan pada Gambar~\ref{fig:logo-lama},
sedangkan logo terbaru ditunjukkan pada Gambar~\ref{fig:logo-baru}.

% ── Contoh gambar: caption DI BAWAH gambar ──
\begin{figure}[h]
  \centering
  \includegraphics[width=3cm]{assets/amikom/contoh/logo_lama.png}
  \caption{Logo lama}\label{fig:logo-lama}
\end{figure}

\begin{figure}[h]
  \centering
  \includegraphics[width=4cm]{assets/amikom/contoh/logo_baru.png}
  \caption{Logo baru}\label{fig:logo-baru}
\end{figure}

Contoh penulisan rumus matematika menggunakan mode matematika
LaTeX. Tingkat akurasi suatu pengujian dapat dihitung menggunakan
persamaan~\eqref{eq:contoh-bab2}.
\begin{equation}
  Akurasi = \frac{TP + TN}{TP + TN + FP + FN} \times 100\%
  \label{eq:contoh-bab2}
\end{equation}
dengan $TP$ adalah \emph{true positive}, $TN$ adalah \emph{true
negative}, $FP$ adalah \emph{false positive}, dan $FN$ adalah
\emph{false negative}.

\section{Penelitian Sebelumnya}
Penelitian sebelumnya berisi rujukan dari jurnal atau prosiding
minimal 3 artikel yang terdiri dari 2 artikel jurnal nasional dan
1 artikel jurnal internasional \parencite{nadzif2019sistem,
haya2020rancang, anuradha2020microcontroller}. Uraikan dengan
jelas kajian pustaka yang menimbulkan gagasan dan mendasari
penelitian yang akan dilakukan, menguraikan teori, temuan, dan
bahan penelitian lain yang dijadikan landasan, serta munculkan
perbedaan Anda sebagai peneliti yang akan mengembangkan sebuah
gagasan, metode, atau hasil yang dicapai. Uraian dalam penelitian
sebelumnya diarahkan untuk menyusun kerangka atau konsep yang akan
digunakan dalam penelitian.

% ── Contoh tabel multi-halaman dengan "(Lanjutan)" ──
% Jika tabel melebihi satu halaman, longtable otomatis meneruskan
% ke halaman berikutnya dengan judul "Tabel X.Y. ... (Lanjutan)"
% dan kepala tabel yang berulang (lihat \TabelLanjutan di
% \endhead). Contoh tabel sengaja dibuat panjang agar menampilkan
% halaman lanjutan.
\begin{tabel}{@{}c l l l@{}}
  \caption{Nilai mata kuliah}\label{tab:nilai}\\
  \hline
  \rowcolor{gray!25}
  No & NIM & Nama Mahasiswa & Nilai\\
  \hline
  \endfirsthead
  \TabelLanjutan{4}{Nilai mata kuliah}\\
  \hline
  \rowcolor{gray!25}
  No & NIM & Nama Mahasiswa & Nilai\\
  \hline
  \endhead
  \hline
  \endfoot
  \endlastfoot
  1 & 11.11.0001 & Mahasiswa 1 & B\\
  2 & 11.11.0002 & Mahasiswa 2 & B\\
  3 & 11.11.0003 & Mahasiswa 3 & A\\
  4 & 11.11.0004 & Mahasiswa 4 & B\\
  5 & 11.11.0005 & Mahasiswa 5 & A\\
  6 & 11.11.0006 & Mahasiswa 6 & C\\
  7 & 11.11.0007 & Mahasiswa 7 & B\\
  8 & 11.11.0008 & Mahasiswa 8 & A\\
  9 & 11.11.0009 & Mahasiswa 9 & B\\
  10 & 11.11.0010 & Mahasiswa 10 & A\\
  11 & 11.11.0011 & Mahasiswa 11 & B\\
  12 & 11.11.0012 & Mahasiswa 12 & C\\
  13 & 11.11.0013 & Mahasiswa 13 & A\\
  14 & 11.11.0014 & Mahasiswa 14 & B\\
  15 & 11.11.0015 & Mahasiswa 15 & A\\
  16 & 11.11.0016 & Mahasiswa 16 & B\\
  17 & 11.11.0017 & Mahasiswa 17 & A\\
  18 & 11.11.0018 & Mahasiswa 18 & B\\
  19 & 11.11.0019 & Mahasiswa 19 & C\\
  20 & 11.11.0020 & Mahasiswa 20 & A\\
  21 & 11.11.0021 & Mahasiswa 21 & B\\
  22 & 11.11.0022 & Mahasiswa 22 & A\\
  23 & 11.11.0023 & Mahasiswa 23 & B\\
  24 & 11.11.0024 & Mahasiswa 24 & A\\
  25 & 11.11.0025 & Mahasiswa 25 & C\\
  26 & 11.11.0026 & Mahasiswa 26 & B\\
  27 & 11.11.0027 & Mahasiswa 27 & A\\
  28 & 11.11.0028 & Mahasiswa 28 & B\\
  29 & 11.11.0029 & Mahasiswa 29 & A\\
  30 & 11.11.0030 & Mahasiswa 30 & B\\
  31 & 11.11.0031 & Mahasiswa 31 & A\\
  32 & 11.11.0032 & Mahasiswa 32 & C\\
  33 & 11.11.0033 & Mahasiswa 33 & B\\
  34 & 11.11.0034 & Mahasiswa 34 & A\\
  35 & 11.11.0035 & Mahasiswa 35 & B\\
  36 & 11.11.0036 & Mahasiswa 36 & A\\
  37 & 11.11.0037 & Mahasiswa 37 & B\\
  38 & 11.11.0038 & Mahasiswa 38 & A\\
  39 & 11.11.0039 & Mahasiswa 39 & B\\
  40 & 11.11.0040 & Mahasiswa 40 & C\\
  41 & 11.11.0041 & Mahasiswa 41 & A\\
  42 & 11.11.0042 & Mahasiswa 42 & B\\
  43 & 11.11.0043 & Mahasiswa 43 & A\\
  44 & 11.11.0044 & Mahasiswa 44 & B\\
  45 & 11.11.0045 & Mahasiswa 45 & A\\
\end{tabel}

% ── Contoh tabel penelitian sebelumnya (minimal 3 artikel:
%    2 jurnal nasional, 1 jurnal internasional) ──
% Tabel ini dibuat LANDSCAPE (halaman miring): halaman PDF
% benar-benar dirotasi saat dilihat/dicetak, dan nomor halaman
% tetap tercetak di kanan atas dalam orientasi miring (sesuai
% contoh halaman landscape di template asli).
\begin{landskape}
\begin{tabel}{@{}c p{4.4cm} p{4.4cm} p{4.4cm} p{3.9cm} p{3.6cm}@{}}
  \caption{Penelitian sebelumnya}\label{tab:penelitian-sebelumnya}\\
  \hline
  \rowcolor{gray!25}
  No & Peneliti & Judul & Hasil Penelitian & Persamaan & Perbedaan\\
  \hline
  \endfirsthead
  \TabelLanjutan{6}{Penelitian sebelumnya}\\
  \hline
  \rowcolor{gray!25}
  No & Peneliti & Judul & Hasil Penelitian & Persamaan & Perbedaan\\
  \hline
  \endhead
  \hline
  \endfoot
  \endlastfoot
  1 & Ahmad Rudiansyah, Mardiono, dan Reza Diharja (2020) &
      Desain Alat Monitoring Kapasitas Tabung Gas LPG 3 Kilogram
      Menggunakan Load Cell Dilengkapi Dengan Deteksi Kebocoran
      Gas Berbasis Internet of Things &
      Monitoring level gas dan deteksi kebocoran gas berbasis
      Internet of Things digunakan untuk memonitoring isi gas
      LPG 3 kg dan kebocoran gas untuk memudahkan pembeli
      memantau level dan memberi peringatan ketika terjadi
      kebocoran gas. &
      Pembuatan alat monitoring level gas dan deteksi kebocoran
      gas berbasis Internet of Things &
      Menggunakan wemos D1 mini sebagai mikrokontroler;\\
  2 & Anisha Fadia Haya, Werman Kasoep, Nefy Puteri Novani (2020) &
      Rancang Bangun Smart Case Sistem Monitoring Tabung Gas
      Elpiji Berbasis Mikrokontroler &
      Mengirim data notifikasi menggunakan GSM ketika terjadi
      kebocoran gas dan memantau informasi isi dengan LCD. &
      Menggunakan load cell sensor sebagai sensor utama yang
      mendeteksi untuk monitoring level gas &
      Mikrokontroler yang digunakan adalah arduino Uno R3;\\
  3 & P. Anuradha, Rajeshwar Rao Arabelli, K. Rajkumar, dan
      J. Ravichander (2020) &
      Microcontroller Based Monitoring and Controlling of LPG
      Leaks Using Internet of Things &
      Pembuatan perangkat yang dapat mendeteksi gas menggunakan
      Arduino Uno ATmega328 dan memantau tingkat kebocoran gas
      melalui ThingSpeak. &
      Sistem untuk mendeteksi kebocoran gas menggunakan
      mikrokontroler &
      Studi ini menggunakan Arduino Uno.\\
\end{tabel}
\end{landskape}

% ── CONTOH HALAMAN LANDSCAPE (opsional) ──────────────────────
% Untuk tabel/gambar lain yang membutuhkan halaman miring,
% bungkus dengan lingkungan landskape seperti contoh di atas:
%
%   \begin{landskape}
%     \begin{tabel}{@{}c p{10cm} p{8cm}@{}}
%       \caption{Judul tabel landscape}\label{tab:landscape}\\
%       \hline
%       No & Kolom A & Kolom B\\
%       \hline
%       \endfirsthead
%       \TabelLanjutan{3}{Judul tabel landscape}\\
%       \hline
%       No & Kolom A & Kolom B\\
%       \hline
%       \endhead
%       \hline
%       \endfoot
%       \endlastfoot
%       1 & Data 1 & Data 2\\
%     \end{tabel}
%   \end{landskape}
%
% Lebar teks saat landscape ≈ 22,7 cm — sesuaikan jumlah & lebar
% kolom agar tidak melebihi batas tersebut. Nomor halaman otomatis
% tampil horizontal di kanan atas orientasi miring.

%============================================================
% BAB III — METODE PENELITIAN
% Struktur: A. Tempat dan Waktu Penelitian (Bila Ada),
% B. Metode Pengumpulan Data, C. Alat dan Bahan Penelitian,
% D. Konsep Penelitian.
%============================================================
\chapter{METODE PENELITIAN}

\section{Tempat dan Waktu Penelitian (Bila Ada)}
% Ganti teks panduan di bawah ini dengan metode Anda.
Skripsi berupa penelitian mandiri merupakan kajian ilmu
pengetahuan atau teknologi di bidang Teknik Informatika.
Penelitian dapat dilaksanakan dengan mempertimbangkan tingkat
kedalaman dan manfaat penelitian tersebut, khususnya untuk
pengayaan ilmu pengetahuan atau penerapan ilmu pengetahuan dan
teknologi.

\section{Metode Pengumpulan Data}
Bagian ini memuat penjelasan secara lengkap dan terinci tentang
cara-cara yang digunakan dalam proses pengumpulan data untuk jenis
data yang diperlukan. Misalnya melalui observasi, wawancara,
eksperimen, atau kuesioner. Jika metode kuesioner digunakan, maka
blangko angket kuesioner harus dilampirkan dalam laporan.

\section{Alat dan Bahan Penelitian}
Berisi sumber data, software dan hardware yang digunakan untuk
melakukan penelitian skripsi.

\section{Konsep Penelitian}
Konsep penelitian ini digunakan sebagai dasar untuk menjawab
pertanyaan-pertanyaan penelitian yang diangkat. Dapat diartikan
sebagai mengalirkan jalan pikiran menurut kerangka logis
(\emph{construct logic}) atau kerangka konseptual yang relevan
untuk menjawab penyebab terjadinya masalah (tidak harus membuat
sistem). Jika membuat produk maka menggunakan metode pengembangan
sistem, jika tidak maka alur disesuaikan dengan penelitian yang
akan dilakukan, misalnya menggunakan metode SPK, algoritma dalam
data mining, framework tatakelola sistem informasi, dan lain-lain.

%============================================================
% BAB IV — HASIL DAN PEMBAHASAN
% Bagian paling penting dari skripsi: memuat semua temuan
% ilmiah yang diperoleh sebagai data hasil penelitian.
% Teks disalin verbatim dari template asli
% (reference/Template_Skripsi.docx).
%============================================================
\chapter{HASIL DAN PEMBAHASAN}

\begin{isi}
Bagian Hasil Penelitian dan Pembahasan merupakan bagian yang
paling penting dari Skripsi karena memuat semua temuan ilmiah
yang diperoleh sebagai data hasil penelitian. Pada bab ini
dipaparkan hasil-hasil dari tahapan penelitian yang dilakukan.
Perhatikan/lihat pada bagian konsep penelitian yang ada di Bab 3.
\end{isi}

%============================================================
% BAB V — PENUTUP
% Struktur: A. Kesimpulan, B. Saran.
% Teks disalin verbatim dari template asli
% (reference/Template_Skripsi.docx).
%============================================================
\chapter{PENUTUP}

\section{Kesimpulan}
\begin{isi}
Kesimpulan merupakan pernyataan singkat, jelas, dan tepat tentang
apa yang diperoleh, memuat keunggulan dan kelemahan, dapat
dibuktikan, serta terkait langsung dengan Rumusan Masalah dan
Tujuan Penelitian. Uraian pada bagian ini harus merupakan
pernyataan yang pernah dianalisis/dibahas pada bagian sebelumnya,
bukan pernyataan yang sama sekali baru dan tidak pernah dibahas
pada bagian sebelumnya, serta merupakan jawaban atas
permasalahan yang dirumuskan. Bagian ini tidak perlu ada uraian
penjelasan lagi.
\end{isi}

\section{Saran}
\begin{isi}
Saran memuat berbagai usulan atau pendapat yang sebaiknya
dikaitkan oleh penelitian sejenis. Saran dibuat berdasarkan
kelemahan, pengalaman, kesulitan, kesalahan, temuan baru yang
belum diteliti dan berbagai kemungkinan arah penelitian
selanjutnya. Saran yang disusun harus didasarkan atas hasil
penelitian yang dilakukan.
\end{isi}


% ── BAGIAN AKHIR (tanpa nomor halaman) ──
%============================================================
% DAFTAR PUSTAKA
%   Bagian ini TIDAK menggunakan nomor halaman.
%   Entri diurutkan alfabetis otomatis (biblatex), format gaya
%   Amikom dikonfigurasi di config/formatting.tex.
%============================================================
\clearpage
\pagestyle{empty}   % mulai bagian tanpa nomor halaman

\chapter*{DAFTAR PUSTAKA}
\thispagestyle{empty}
\BagianAkhir{DAFTAR PUSTAKA}

% Entri yang belum disitasi tetap dicetak (contoh template).
% Hapus baris berikut saat sudah menulis skripsi sungguhan.
\nocite{nadia2019tantangan, albarda2004strategi, wibisono2013menentukan}

\printbibliography[heading=none]

%============================================================
% LAMPIRAN
%   Berisi kartu bimbingan, surat keterangan instansi, blangko
%   kuesioner, instrument penelitian, listing program, dll.
%   Kartu bimbingan WAJIB dimasukkan. Judul lampiran:
%   "Lampiran 1. Kartu Bimbingan Dosen Pembimbing 1".
%   Bagian ini TIDAK menggunakan nomor halaman.
%============================================================
\chapter*{LAMPIRAN}
\thispagestyle{empty}
\BagianAkhir{LAMPIRAN}

\noindent
Lampiran berisi kartu bimbingan, surat keterangan dari instansi
tempat penelitian, blangko kuesioner, instrumen penelitian,
listing program, peraturan-peraturan, dan sebagainya yang
berfungsi melengkapi laporan penelitian. Kartu bimbingan wajib
dimasukkan ke lampiran.

% ── Lampiran 1 ──
\Lampiran{Kartu Bimbingan Dosen Pembimbing 1}
% Masukkan hasil scan kartu bimbingan yang sudah ditandatangani.
% Cara memasukkan scan PDF (gunakan paket pdfpages):
%   \includepdf[pages=1]{assets/kartu-bimbingan-1.pdf}
\KotakScan{Tempat hasil scan Kartu Bimbingan Dosen Pembimbing 1\\
yang sudah ditandatangani}

% ── Lampiran 2 ──
\Lampiran{Kartu Bimbingan Dosen Pembimbing 2}
%   \includepdf[pages=1]{assets/kartu-bimbingan-2.pdf}
\KotakScan{Tempat hasil scan Kartu Bimbingan Dosen Pembimbing 2\\
yang sudah ditandatangani}

% ── Lampiran 3 dan seterusnya (contoh: blangko kuesioner) ──
% Tambahkan \Lampiran{...} sesuai kebutuhan penelitian Anda.
\Lampiran{Blangko Kuesioner}
% Jika metode kuesioner digunakan, blangko angket kuesioner
% harus dilampirkan dalam laporan.
\KotakScan{Tempat blangko kuesioner penelitian}


\end{document}
